% Generated draft: requires mathematical and source review.
\documentclass[11pt]{article}
\usepackage[margin=25mm]{geometry}
\usepackage{fontspec}
\setmainfont{Latin Modern Roman}
\usepackage{amsmath,amssymb,mathtools,unicode-math}
\setmathfont{Latin Modern Math}
\usepackage{enumitem,needspace,xurl,hyperref}
\setlength{\parindent}{0pt}
\setlength{\parskip}{5pt}
\setlength{\emergencystretch}{3em}
\newcommand{\N}{\mathbb{N}}
\newcommand{\Z}{\mathbb{Z}}
\newcommand{\Q}{\mathbb{Q}}
\newcommand{\R}{\mathbb{R}}
\newcommand{\C}{\mathbb{C}}
\newcommand{\F}{\mathbb{F}}
\newcommand{\A}{\mathbb{A}}
\newcommand{\PP}{\mathbb{P}}
\newcommand{\E}{\mathbb{E}}
\newcommand{\eps}{\varepsilon}
\newcommand{\dd}{\,\mathrm{d}}
\begin{document}
\section*{UnsolvedMath ID: 30001156}
\textbf{Source problem code:} catalogue entry
\subsection{Definitions and mathematical statement}
Let \(\Omega\subset\R^2\) be a nonempty open convex domain with finite positive area \(A=|\Omega|\) and finite perimeter \(L=\mathcal H^1(\partial\Omega)\). Interpret \(\lambda_1(\Omega)\) as the first Dirichlet eigenvalue of the Laplacian, namely
\[
\lambda_1(\Omega)=\inf_{0\ne u\in H_0^1(\Omega)}\frac{\int_\Omega |\nabla u|^2\,\dd x}{\int_\Omega |u|^2\,\dd x},
\]
where \(H_0^1(\Omega)\) is the closure of \(C_c^\infty(\Omega)\) in the Sobolev space \(H^1(\Omega)\). Let \(j_{0,1}\) denote the first positive zero of the Bessel function \(J_0\) of the first kind.

The question is whether
\[
\exists K>0\quad\forall\Omega\text{ as above},\qquad
\lambda_1(\Omega)\geq\frac{\pi j_{0,1}^2}{A}+K\frac{L^2-4\pi A}{A^2}.
\]
The constant \(K\) must be independent of \(\Omega\). The supplied original statement attributes to the report the contrary conjecture: for every \(K>0\), there exists such a convex domain \(\Omega\) satisfying
\[
\lambda_1(\Omega)<\frac{\pi j_{0,1}^2}{A}+K\frac{L^2-4\pi A}{A^2}.
\]
In the terminology of the supplied record, this asks whether the optimal coefficient of the isoperimetric defect \(L^2-4\pi A\) is zero.

The Dirichlet boundary condition and the definitions of \(\lambda_1\) and \(j_{0,1}\) are the standard interpretation suggested by the supplied Faber--Krahn context; they are not explicitly specified in the supplied statement and require source confirmation. The domain conventions above make the quantities precise. The original statement begins with a simply connected planar domain but explicitly restricts the proposed inequality to convex domains; convex domains are simply connected, so the question here retains the convex restriction.
\subsection{Short English statement}
Does the first Dirichlet eigenvalue of every planar convex domain exceed the disk's equal-area eigenvalue by a universal positive multiple of its perimeter defect divided by the square of its area? The supplied original record attributes the conjecture that no positive coefficient works to the report. The boundary-condition interpretation requires confirmation, and the supplied open label is not independently verified.
\subsection{Sources}
\begin{itemize}
\item Local UnsolvedMath registry: supplied complete record for permanent ID \(30001156\), source problem code \texttt{catalogue entry}, titled \emph{Isoperimetric-Defect Corrections to Faber--Krahn Bounds}, in the Oberwolfach Reports Open Problems collection. The proposed year is recorded as 2009; no individual proposer is identified.
\item Supplied primary reference: \emph{Low Eigenvalues of Laplace and Schr\"odinger Operators} (2009), Oberwolfach Reports, pp.\ 355--428, DOI \texttt{10.4171/owr/2009/06}. \url{https://doi.org/10.4171/owr/2009/06}. The supplied original statement attributes the conjecture that no positive \(K\) exists to this report; no exact problem page or named speaker is supplied.
\item Supplied literature reference, labeled ``quantitative isoperimetric survey'': \url{https://doi.org/10.1007/s13373-015-0074-x}. No authors or full bibliographic title are supplied for this reference.
\item Supplied literature assessment checked on \texttt{2026-08-21}: status recorded as open; the curation reports finding no proof or counterexample for this exact positive-correction bound on convex domains and states that available quantitative Faber--Krahn results concern different deficits and do not settle this perimeter-defect formulation. This is supplied dated curation, not an independently verified literature assessment.
\item The registry labels the statement \texttt{corrected\_verified} and reports that its mathematical content was checked against the official report, with notation, context, and question form normalized without changing the source-level claim. This verification claim is supplied metadata and has not been independently checked for this draft.
\end{itemize}
\end{document}
